\documentstyle[% 


righttag% 


,11pt% 


,amssymb% 


,russian%               remove in Eng. 


% ,izvru%                 Set izven% in Eng. 


,izvru-ks%


]{amsart} 


 


% Authorized


\newcommand{\grad}{\operatorname{grad}}


\newcommand{\Hess}{\operatorname{Hess}}





\theoremstyle{definition} 


\theorembodyfont{\rm} 


\newtheorem{examnonum}{\hskip\parindent Пример} 


\renewcommand{\theexamnonum}{} 


\renewcommand{\baselinestretch}{1.05}


 


%\hoffset=-15mm%                  For Eng. 


 


\setcounter{page}{96} 


\god{1998} 


\nomer{6~(433)} %                Remove in Eng. 


%\nomer{Vol.\,42, No.\,6}%        Set in Eng. 


%\str{\pageref{begin}--\pageref{end}}%                    Set in Eng. 


\udk{514.75} 


\author{\it М.А.\,Чешкова} 


% NB:   ^^ remove in Eng.


\title{К геометрии центральной проекции


$\size{14}{18pt}\selectfont n$-поверхности\\


в евклидовом пространстве $E^{n+m}$ }


 


%\thanks{Работа выполнена при поддержке РФФИ, грант \No. 94-01-00183.} 


\date{} 


% \pagestyle{myheadings}%         Set in Eng. 


 


\pagestyle{plain} %              Remove in Eng. 


\begin{document} 


 


% \thispagestyle{plain}%          Set in Eng. 


 


\maketitle 


\label{begin} 


В евклидовом пространстве $E^{n+m}$ рассматриваются две гладкие 


$n$-поверхности $M$, $\overline {M}$ и диффеоморфизм


$f:M \rightarrow  \overline{M}$.


Исследуется случай, когда $f$ --- центральная проекция.                                 





Обозначим через $r$ радиус-вектор точки $p \in M$, через $\overline{r}$


--- радиус-вектор точки $f(p) \in {\overline M}$. 


Тогда отображение $f: M \rightarrow \overline{M}$  запишется в виде


$\overline{r}=lr$, $l\in F(M).$


Дифференциал отображения $f$   определится из равенства


$df(X)= df(\partial _{X} r) = \partial _{X} \overline{r}$, $X \in TM. $


Имеем


$ df(X)=(Xl)r+lX$, $X \in TM. $


Отображение $f$  индуцирует на $M$  метрику


$$


\overline{g}(X,Y)=\langle df(X),df(Y)\rangle =


l^2 g(X,Y)+ (Xl)(Yl)\langle


r,r\rangle  + l(Yl)\langle X,r\rangle+ l(Xl)\langle Y,r\rangle.


$$


Положим


$r = U+\tau$, $U \in TM$, $\tau \in TM^{\bot}$,


$\langle r,r\rangle=2\rho .$


Имеем


\begin{equation}


\overline{g}(X,Y)= l^{2}g(X,Y) +\Psi(X,Y),


\end{equation}


где


$\Psi(X,Y)=\gamma (X)\omega(Y)+\gamma(Y)\omega(X)$,


$\gamma(X)=Xl$, $\omega(X)=\rho(Xl)+l\langle U,X\rangle.$





\begin{lem}


Cледующие утверждения эквивалентны\,{\em :}





$1)$ центральная проекция $f: M \rightarrow \overline{M}$


--- конформное отображение\,$;$





$2)$ $\Psi =0$.


\end{lem}





\begin{pf} Отображение $f: M \rightarrow \overline{M}$ конформно


тогда и только тогда, когда $\overline{g}(X,Y)= \sigma^{2}g(X,Y)$,


$\sigma \in F(M),$ или $(\sigma^2 -l^2 )g(X,Y)= \Psi(X,Y)$.


А так как ранг квадратичной формы $\Psi <3 $, то при $n>2$\;


$\sigma^2 -l^2 =0$, $\Psi =0$.


Для $n=2$ существуют изотермические координаты $u$, $v$. Тогда


$\gamma_{u}\omega_{v}+\gamma_{v}\omega_{u}=0$, $\gamma_{v}\omega_{v}=


\gamma_{u}\omega_{u}$, откуда следует $\Psi =0$.


\end{pf}





\begin{cor}                   


Если  центральная проекция есть конформное отображение, то


либо $f$ есть гомотетия, либо $V = -\frac{1}{\rho}U = \grad(\ln|l|)$.


\end{cor}





\begin{pf} $\Psi =0$ в случае, когда


$\gamma = 0$, $l=\text{const}$, $f$ --- гомотетия, либо


$\omega = 0$, $X \ln|l|+ \frac{1}{\rho}\langle U,X\rangle =0,$ т.\,е.


вектор $V = - \frac{1}{\rho} U$  есть  $\grad \ln|l|$.


\end{pf}





\begin{cor}


Если  центральная проекция $f: M \rightarrow \overline{M}$  есть


конформное отображение, то $f$ --- либо гомотетия, либо --- инверсия.


\end{cor}





\begin{pf}


Если $l=\text{const}$, то $f$ --- гомотетия, если $l\ne\text{const}$,


то $X \ln|l| = -\frac{1}{\rho}\langle X,U\rangle$. 


С другой стороны, дифференцируя


равенство $2\rho =\langle r,r\rangle$, получим


$X\rho =\langle X,r\rangle=\langle X,U\rangle$. Откуда


$l\rho= k=\text{const}$, $\overline{r}= \frac{2kr}{\langle r,r\rangle}$,


т.\,е. $f$ --- инверсия. 


\end{pf}





Пусть $\overline{\alpha}$ --- вторая фундаментальная форма $n$-поверхности


$\overline{M}$, $\overline{\nabla}$ --- связность  Леви-Чивита метрики


$\overline{g}$. Имеет место  [1] равенство


\begin{equation}


\partial _{X}df Y-df\overline{\nabla}_{X}Y=  \overline{\alpha}(X,Y),


\end{equation} 


откуда


$\partial_{X}((Yl)r+lY)-(\overline{\nabla}_{X}Y l)r-


l(\overline{\nabla}_{X}Y)=


\overline{\alpha}(X,Y)$.


Используя уравнения Гаусса-Вейнгартена ([2], с.\,23), получим


$$


(\Hess_{X,Y}^{\overline{\nabla}}l)r+(Xl)Y+(Yl)X+


l(\nabla_{X}Y-\overline{\nabla}_{X}Y)+


l\alpha(X,Y)=\overline{\alpha}(X,Y),


$$


где $\Hess_{X,Y}^{\overline{\nabla}}l=XYl-\overline{\nabla}_{X}Y l$


--- гессиан функции $l$ в связности $\overline{\nabla}$.


Разлагая $\overline{\alpha}$ на касательную


$\overline{\alpha}^{\top}$


и нормальную $\overline{\alpha}^{\bot}$ cоставляющие, получим


\begin{gather}


\overline{\alpha}^{\top}(X,Y)=(Xl)Y+(Yl)X+


l(\nabla_{X}Y-\overline{\nabla}_{X}Y)+


(\Hess_{X,Y}^{\overline{\nabla}}l)U, \\


%


\overline{\alpha}^{\bot}(X,Y)=


l\alpha(X,Y)+ (\Hess_{X,Y}^{\overline{\nabla}}l)\tau.\notag


\end{gather}


Векторы $\overline{\alpha}(X,Y)$ определяют первое нормальное пространство


$\overline{N}$ или главную нормаль поверхности $\overline{M}$.


Из (3) вытекают следующие теоремы.    





\begin{thm}


Если $f: M \rightarrow \overline{M}$  --- центральная проекция, то следующие


утверждения эквивалентны\,{\em :}





$1)$ главная нормаль $\overline{N}$ поверхноcти $\overline{M}$


параллельна поверхности $M$,





$2)$ $l\alpha(X,Y)=-(\Hess_{X,Y}^{\overline{\nabla}}l)\tau.$


\end{thm}





\begin{cor}                   


Если $f: M \rightarrow \overline{M}$  --- центральная проекция, и 


главная нормаль $\overline{N}$ поверхности $\overline{M}$ параллельна


поверхности


$M$, то главная нормаль поверхности $M$ одномерна и параллельна полю $\tau$.


\end{cor}





\begin{thm}


Если $f: M \rightarrow \overline{M}$  --- центральная проекция, то следующие


утверждения эквивалентны\,{\em :}





$1)$ главная нормаль поверхности $\overline{M}$ ортогональна поверхности $M$,





$2)$ $\overline{\nabla}_{X}Y= \nabla_{X}Y+(X\ln l)Y+ (Y\ln l)X+


\frac{1}{l}(\Hess_{X,Y}^{\overline{\nabla}}l)U$.


\end{thm}  





Исследуем $\overline{\alpha}$, когда $f: M \rightarrow \overline{M}$


--- конформное отображение.


Если $f$ --- гомотетия, то $Xl=0$, $\nabla=\overline{\nabla}$,


$\overline{\alpha}^{\top}=0$, $\overline{\alpha}^{\bot}(X,Y)=l\alpha(X,Y).$                    


Если  $f: M \rightarrow \overline{M}$ --- конформное отображение, 


не являющееся гомотетией, то 


\begin{gather*}


\overline{\nabla}_{X}Y=\nabla_{X}Y+(X\ln|l|)Y+(Y\ln|l|)X-


\frac{\rho}{l}(\Hess_{X,Y}^{\overline{\nabla}} l)V-\frac{1}{l}


\overline{\alpha}^{\top}(X,Y), \\


%


\overline{g}=l^{2}g,\quad V=\grad(\ln|l|).


\end{gather*}


С другой стороны, если $f$ --- конформное отображение, то ([3], с.\,18)


$\overline{\nabla}_{X}Y=\nabla_{X}Y+(X\ln|l|)Y+(Y\ln|l|)X-g(X,Y)V,$


где $\overline{g}=l^{2}g$, $V=\grad(\ln|l|).$


Таким образом, имеем                            


\begin{equation}


\overline{\alpha}^{\top}(X,Y)=


(l g(X,Y)+\rho Hess_{X,Y}^{\overline{\nabla}})V.


\end{equation}





Из (3), (4) вытекают





\begin{thm}


Если центральная проекция $f$ есть конформное отображение, то


$\overline{\alpha} (X,Y)= l\alpha (X,Y)+ (\Hess_{X,Y}^{\overline{\nabla}})r +


l g(X,Y)V.$


\end{thm}





\begin{cor}


Если главная нормаль поверхности $\overline{M}$ ортогональна поверхности $M$,


и центральная проекция есть конформное отображение, отличное от гомотетии, 


то


$\overline{\alpha} ^{\bot}(X,Y)= \overline{\alpha}(X,Y)=l\alpha(X,Y)+


\frac{l}{\rho} g(X,Y)\tau$.


\end{cor} 





\begin{examnonum} Центральная проекция $f$ есть конформное отображение на


$n$-плоскость. Тогда $\alpha(X,Y)=-\frac{1}{\rho}g(X,Y) \tau.$


Положим $\nu$ --- орт $\tau$. Тогда $\alpha(X,Y)=kg(X,Y)\nu$, 


где $k=\frac{1}{\rho|\tau|}$, $A_{\nu}X=kX.$


Из уравнений Гаусса-Кодацци ([2], с.\,29) следует


$(Xk)\nu+k\nabla^{\bot}_{X}\nu=0$. Откуда имеем


$Xk=0$, $\nabla^{\bot}_{X}\nu=0$, т.\,е. поверхность $M$ принадлежит


гиперсфере с центром $C=r+\frac{1}{k}\nu$ радиуса $\frac{1}{k}$ и 


$(n+1)$-плоскости $(\overline{M},O)$, т.\,е. является  $n$-сферой.


\end{examnonum}





\begin{thebibliography}{9}


\bibitem{1}


Чешкова М.А. {\em К геометрии $n$-поверхностей в евклидовом 


пространстве $E^{2n+1}$} // Дифференц. геометрия многообразий фигур. --


1997. -- Вып.\,28. -- С.78--81.


\bibitem{2}


Кобаяси Ш., Номидзу К. {\em Основы дифференциальной геометрии}. -- T.\,2. --


М.: Наука, 1981. -- 414 с.


\bibitem{3}


Chen B.-Y. {\em Geometry of submanifolds and its applications}. --


Tokyo: Sci. Univ., 1981. -- 96~p.





\end{thebibliography}





\vskip 2cm





\postup{\it Алтайский государственный}{$\qquad$\it Поступили}


\postup{\it университет}{{\it полный текст} 16.12.1996} 


\postup{}{{\it краткое собщение} 27.01.1998} 


\label{end} 


\end{document} 


